\clearpage
\SourceLeaf{404}
\begin{center}\large [VERALLGEMEINERUNG DES SATZES VON LEGENDRE.]\end{center}
\begin{center}[I.]\end{center}
\begin{center}\textsc{Gauss} an \textsc{Olbers}. Göttingen, 9. Oktober 1825.\end{center}
\indent\mbox{. . . . . . . . . .} Ich habe in diesen Tagen mit Blick auf mein künftiges Werk über Höhere Geodäsie angefangen, einen (sehr) kleinen Teil des Stoffs über die krummen Flächen in Gedanken etwas zu ordnen. Ich komme aber zu der Überzeugung, dass mich die Eigenart meiner gesamten Behandlung zwingt, um des Zusammenhangs willen \textls{sehr weit auszuholen}. Ich muss sogar meine Auffassung von den Krümmungsradien bei \textls{ebenen} Kurven voranstellen. Daher bin ich fast unsicher geworden, ob es nicht ratsamer wäre, einen Teil dieser Lehren, der rein geometrisch (in analytischer Form) ist und Neues mit Bekanntem in neuer Form verbindet, zunächst gesondert auszuarbeiten, ihn vielleicht aus dem Werk herauszulösen und als eine oder zwei Abhandlungen in unsere Commentationen aufzunehmen. Vorerst kann ich die Form der Veröffentlichung aber offenlassen und damit fortfahren, alles zu Papier zu bringen. \mbox{. . .}
